% Generated draft: requires mathematical and source review.
\documentclass[11pt]{article}
\usepackage[margin=25mm]{geometry}
\usepackage{fontspec}
\setmainfont{Latin Modern Roman}
\usepackage{amsmath,amssymb,mathtools,unicode-math}
\setmathfont{Latin Modern Math}
\usepackage{enumitem,needspace,xurl,hyperref}
\setlength{\parindent}{0pt}
\setlength{\parskip}{5pt}
\setlength{\emergencystretch}{3em}
\newcommand{\N}{\mathbb{N}}
\newcommand{\Z}{\mathbb{Z}}
\newcommand{\Q}{\mathbb{Q}}
\newcommand{\R}{\mathbb{R}}
\newcommand{\C}{\mathbb{C}}
\newcommand{\F}{\mathbb{F}}
\newcommand{\A}{\mathbb{A}}
\newcommand{\PP}{\mathbb{P}}
\newcommand{\E}{\mathbb{E}}
\newcommand{\eps}{\varepsilon}
\newcommand{\dd}{\,\mathrm{d}}
\begin{document}
\section*{UnsolvedMath ID: 30004721}
\textbf{Source problem code:} catalogue entry
\subsection{Definitions and mathematical statement}
For a smooth manifold \(M\), under the convention that no completeness or boundary conditions are imposed, let
\[
\mathcal R_{>0}(M)=\{g:\ g\text{ is a smooth Riemannian metric on }M,\ \operatorname{Scal}_g(x)>0\text{ for every }x\in M\}.
\]
Here \(\operatorname{Scal}_g\) denotes scalar curvature. The source does not specify whether metrics must be complete, whether manifolds are closed, or what conditions apply at a boundary; thus this convention is provisional.

To express the question for ordinary framed interior surgery, let \(M\) have dimension \(n\), let \(0\leq p\leq n-1\), and let
\[
\iota:S^p\times D^{n-p}\hookrightarrow\operatorname{int}(M)
\]
be a smooth embedding specifying a framed tubular neighborhood. Here \(S^k\) and \(D^k\) denote the standard sphere and closed disk. The manifold obtained by this surgery is
\[
M_\iota=
\left(M\setminus\iota\bigl(S^p\times\operatorname{int}(D^{n-p})\bigr)\right)
\mathop{\cup}_{S^p\times S^{n-p-1}}
\left(D^{p+1}\times S^{n-p-1}\right),
\]
where the gluing uses the boundary identification specified by \(\iota\), with the usual smoothing. The surgery sphere has codimension \(n-p\).

\textbf{Problem.} Determine which kinds of surgery preserve nonexistence of positive scalar curvature metrics: find conditions on the dimension, surgery sphere, framing, embedding, and any necessary manifold hypotheses that ensure
\[
\mathcal R_{>0}(M)=\varnothing
\quad\Longrightarrow\quad
\mathcal R_{>0}(M_\iota)=\varnothing
\]
for every manifold and every surgery satisfying those conditions. Ordinary framed interior surgery gives one standard interpretation of the source question; the source does not specify the intended range of surgeries or manifold hypotheses. It therefore supplies a broad classification question rather than a single fully specified universal conjecture.
\subsection{Short English statement}
Determine which surgeries preserve the property that a smooth manifold admits no Riemannian metric with scalar curvature strictly positive everywhere. The source leaves the dimension, manifold class, metric completeness, boundary conditions, and permitted surgeries unspecified, so these require clarification before a definitive formulation can be fixed.
\subsection{Sources}
\begin{itemize}
\item Local UnsolvedMath registry, supplied record with permanent ID \(30004721\), source problem code \texttt{catalogue entry}, titled ``Preservation of Positive-Scalar-Curvature Obstructions Under Surgery.'' The record belongs to the Oberwolfach Reports Open Problems collection and gives proposed year \(2021\); no individual proposer is named.
\item Supplied report citation: \emph{Analysis, Geometry and Topology of Positive Scalar Curvature Metrics} (2022), Oberwolfach Reports, pp.\ 1579--1649, DOI \texttt{10.4171/owr/2021/30}. Supplied source URL and literature-source link, labeled ``Oberwolfach Reports 2022'': \url{https://doi.org/10.4171/owr/2021/30}. The background describes these reports as documenting mathematical workshops and their open-problem discussions.
\item Literature assessment dated exactly \texttt{2026-08-22}: the supplied curation labels the problem open and reports that the forward codimension-three surgery theorem preserves existence of positive scalar curvature, whereas no general reverse or nonexistence surgery criterion answers the broad question stated here. This is supplied dated literature triage derived from the OWR curation, not an independently verified assessment; time-sensitive status and cited claims require independent verification.
\item The supplied statement status is \texttt{corrected\_verified}. Its accompanying explanation describes a publication-format audit that preserved mathematical content while repairing notation or typography; this is recorded provenance, not mathematical verification of this draft.
\end{itemize}
\end{document}
